\begin{figure}[tb]
\centering
\begin{tikzpicture}[x={(1cm,0cm)},y={(.35cm,.35cm)},z={(0cm,1cm)},scale=1.05]
 \draw[coordinate axis] (-3,0,0)--(3.2,0,0) node[right]{$i_d$};
 \draw[coordinate axis] (0,-2,0)--(0,2.1,0) node[above right]{$i_q$};
 \draw[coordinate axis] (0,0,-2)--(0,0,2.25) node[above]{$i_e$};
 \begin{scope}[canvas is xy plane at z=-1.25]
  \draw[excitation slice,estimated quantity] (.8,0) ellipse (1.35 and .72);
 \end{scope}
 \begin{scope}[canvas is xy plane at z=0]
  \draw[excitation slice,derived quantity] (0,0) ellipse (1.35 and .72);
 \end{scope}
 \begin{scope}[canvas is xy plane at z=1.25]
  \draw[excitation slice,reference quantity] (-.8,0) ellipse (1.35 and .72);
 \end{scope}
 \draw[projection line] (.8,0,-1.25)--(0,0,0)--(-.8,0,1.25);
 \node[annotation] at (2.25,0,-1.25){$i_e<0$};
 \node[annotation] at (1.7,0,0){$i_e=0$};
 \node[annotation] at (.85,0,1.25){$i_e>0$};
\end{tikzpicture}
\caption{Constant-excitation cuts are congruent ellipses.  Their centers move
on the null line, generating the cylinder.}
\label{fig:excitation-slices}
\end{figure}
